% Aula 02 --- derivar as fórmulas do Ridge e do Lasso no caso ortonormal.
As notas da aula de Regressão Linear apresentam duas fórmulas prontas para o caso ortonormal,
$\mathbb{X}^\top\mathbb{X}=\mathbb{I}$, e não as derivam. São elas que explicam a
frase "o Ridge encolhe e o Lasso zera", e o exercício é preencher essa lacuna.

Use as definições da aula, \textbf{sem} fator $\tfrac12$ no RSS:
\[
  \widehat{\bbeta}^{\,\text{ridge}}
   = \argmin_{\bbeta}\Big\{\norm{\bm y-\mathbb{X}\bbeta}^2+\lambda\sum_j\beta_j^2\Big\},
  \qquad
  \widehat{\bbeta}^{\,\text{lasso}}
   = \argmin_{\bbeta}\Big\{\norm{\bm y-\mathbb{X}\bbeta}^2+\lambda\sum_j\abs{\beta_j}\Big\}.
\]
Aqui $\mathbb{X}\in\R^{n\times p}$ não tem a coluna de uns: não há intercepto, e as
duas penalidades somam sobre os $p$ coeficientes.
\begin{itens}
  \item \pts{1,1} Comece pelo passo que torna tudo possível. Supondo
        $\mathbb{X}^\top\mathbb{X}=\mathbb{I}$, mostre que o estimador de mínimos
        quadrados é $\widehat\bbeta=\mathbb{X}^\top\bm y$ e que
        \[
          \norm{\bm y-\mathbb{X}\bbeta}^2
            = \norm{\bm y}^2-\norm{\widehat\bbeta}^2
              + \sum_{j=1}^{p}\big(\beta_j-\widehat\beta_j\big)^2 .
        \]
        Explique por que essa identidade faz o problema \textbf{se separar} em $p$
        problemas de uma variável --- e por que é justamente isso que falha quando
        as colunas não são ortonormais.
  \item \pts{0,7} Com a separação em mãos, derive
        $\widehat\beta_j^{\,\text{ridge}}=\widehat\beta_j/(1+\lambda)$.
  \item \pts{1,4} Agora o Lasso. A dificuldade nova é que $\abs{\beta_j}$ não é
        derivável em zero, então não dá para igualar a derivada a zero e pronto:
        trate separadamente as regiões $\beta_j>0$, $\beta_j<0$ e o ponto
        $\beta_j=0$. Derive
        $\widehat\beta_j^{\,\text{lasso}}
          =\operatorname{sinal}(\widehat\beta_j)\big(\abs{\widehat\beta_j}-\lambda/2\big)_+$
        e diga qual é a condição exata que leva um coeficiente a zero.
  \item \pts{0,8} O limiar deu $\lambda/2$, e não $\lambda$. De onde veio esse
        $2$? Diga o que mudaria nas duas fórmulas se a aula tivesse definido o RSS
        com um $\tfrac12$ na frente, e que lição prática isso deixa.
\end{itens}

\begin{solucao}
\textbf{(a)} Com $\mathbb{X}^\top\mathbb{X}=\mathbb{I}$ a matriz é invertível, e as
equações normais $\mathbb{X}^\top\mathbb{X}\,\bbeta=\mathbb{X}^\top\bm y$ dão
diretamente $\widehat\bbeta=\mathbb{X}^\top\bm y$.

Para a identidade, o caminho mais curto é o teorema de Pitágoras. Escreva
\[
  \bm y-\mathbb{X}\bbeta=(\bm y-\mathbb{X}\widehat\bbeta)+\mathbb{X}(\widehat\bbeta-\bbeta).
\]
As duas parcelas são ortogonais, porque as equações normais dizem exatamente que o
resíduo do MQO é ortogonal às colunas de $\mathbb{X}$:
$\mathbb{X}^\top(\bm y-\mathbb{X}\widehat\bbeta)=\bm 0$. Logo, para \emph{qualquer}
$\mathbb{X}$ de posto cheio,
\[
  \norm{\bm y-\mathbb{X}\bbeta}^2
  =\norm{\bm y-\mathbb{X}\widehat\bbeta}^2
   +(\bbeta-\widehat\bbeta)^\top\mathbb{X}^\top\mathbb{X}\,(\bbeta-\widehat\bbeta).
  \tag{$*$}
\]
No caso ortonormal a forma quadrática é $\sum_j(\beta_j-\widehat\beta_j)^2$, e
\[
  \norm{\bm y-\mathbb{X}\widehat\bbeta}^2
  =\norm{\bm y}^2-2\,\widehat\bbeta^\top\underbrace{\mathbb{X}^\top\bm y}_{=\;\widehat\bbeta}
   +\widehat\bbeta^\top\underbrace{\mathbb{X}^\top\mathbb{X}}_{=\;\mathbb{I}}\widehat\bbeta
  =\norm{\bm y}^2-\norm{\widehat\bbeta}^2 ,
\]
o que dá a identidade do enunciado.

\emph{Por que ela separa o problema.} Somando a penalidade, os dois objetivos ficam
\[
  \norm{\bm y}^2-\norm{\widehat\bbeta}^2+\sum_{j=1}^p h_j(\beta_j),
  \qquad
  h_j(\beta)=(\beta-\widehat\beta_j)^2+\lambda\,\mathrm{pen}(\beta),
\]
com $\mathrm{pen}(\beta)=\beta^2$ no Ridge e $\mathrm{pen}(\beta)=\abs{\beta}$ no
Lasso. A constante não muda o minimizador, e o resto é uma soma de funções de
\emph{uma coordenada cada}. Para somas assim vale um fato simples: $\bbeta^\ast$
minimiza $\sum_jh_j(\beta_j)$ se e só se cada $\beta_j^\ast$ minimiza a sua $h_j$. (Se
cada parcela está no seu mínimo, a soma está; e se alguma não estivesse, trocar só
aquela coordenada baixaria a soma.) Um problema em $\R^p$ virou $p$ problemas na reta.

\emph{Por que isso falha sem ortonormalidade.} Em $(*)$, a forma quadrática tem os
termos cruzados $2(\mathbb{X}^\top\mathbb{X})_{j\ell}(\beta_j-\widehat\beta_j)(\beta_\ell-\widehat\beta_\ell)$,
e o melhor valor de uma coordenada passa a depender das outras. É por isso que o Lasso
não tem fórmula fechada no caso geral, e as fórmulas deste exercício pedem uma
hipótese forte. (É também por isso que o enunciado tira o intercepto: a coluna de uns
tem norma $\sqrt n$ e não caberia numa matriz com $\mathbb{X}^\top\mathbb{X}=\mathbb{I}$.
Com dados centrados o intercepto sai de cena, e é nesse cenário que o caso ortonormal
faz sentido.)

\smallskip
\textbf{(b)} Para o Ridge, $h(\beta)=(\beta-\widehat\beta_j)^2+\lambda\beta^2$ é uma
parábola com $h''=2(1+\lambda)>0$: estritamente convexa, com mínimo único onde a
derivada se anula,
\[
  h'(\beta)=2(\beta-\widehat\beta_j)+2\lambda\beta=0
  \iff
  \widehat\beta_j^{\,\text{ridge}}=\frac{\widehat\beta_j}{1+\lambda}.
\]
O Ridge \textbf{multiplica} todo coeficiente pelo mesmo fator $1/(1+\lambda)$: encolhe
proporcionalmente --- quem é grande perde mais em valor absoluto --- e nunca zera um
coeficiente não nulo, para $\lambda$ finito.

\smallskip
\textbf{(c)} Agora $h(\beta)=(\beta-b)^2+\lambda\abs{\beta}$, com $b=\widehat\beta_j$ e
$\lambda>0$. A função é contínua em toda a reta e derivável fora do zero:
\[
  h'(\beta)=2(\beta-b)+\lambda\quad(\beta>0),
  \qquad
  h'(\beta)=2(\beta-b)-\lambda\quad(\beta<0).
\]
A estratégia é ver onde $h$ cresce e onde decresce, dos dois lados do zero.

\emph{Caso $b>\lambda/2$.} Para $\beta<0$, $h'(\beta)=2\beta-(2b+\lambda)<0$: $h$
decresce em toda a semirreta negativa. Para $\beta>0$, $h'$ troca de sinal uma única
vez, em $\beta=b-\lambda/2>0$: $h$ decresce até ali e cresce depois. Como $h$ é
contínua no zero, o mínimo global é $\beta=b-\lambda/2$.

\emph{Caso $b<-\lambda/2$.} Simétrico: o mínimo é $\beta=b+\lambda/2<0$.

\emph{Caso $\abs{b}\le\lambda/2$.} Para $\beta>0$, $h'(\beta)=2\beta+(\lambda-2b)>0$,
porque $\lambda-2b\ge0$; para $\beta<0$, $h'(\beta)=2\beta-(\lambda+2b)<0$, porque
$\lambda+2b\ge0$. A função decresce até o zero e cresce depois dele: o mínimo é
$\beta=0$ --- justamente o ponto em que ela não é derivável, e que nenhuma equação
"$h'=0$" encontraria.

Os três casos cabem numa fórmula só,
\[
  \widehat\beta_j^{\,\text{lasso}}
  =\operatorname{sinal}(\widehat\beta_j)\big(\abs{\widehat\beta_j}-\lambda/2\big)_+ ,
  \qquad\text{e o coeficiente zera}\iff
  \boxed{\;\abs{\widehat\beta_j}\le\lambda/2\;}.
\]
(Para quem conhece subgradientes: $0$ minimiza a função convexa $h$ quando
$0\in\partial h(0)=-2b+\lambda[-1,1]$, que é a mesma condição.)

Compare com (b): o Lasso \textbf{subtrai} a mesma quantidade, $\lambda/2$, de todo
coeficiente, em vez de multiplicar por um fator. Quem tinha módulo menor que isso não
sobrevive ao deslocamento e vira exatamente zero. É o \emph{soft thresholding}, e é a
razão algébrica de o Lasso fazer seleção de variáveis.

\smallskip
\textbf{(d)} O $2$ vem de derivar o quadrado. A derivada de $(\beta-\widehat\beta_j)^2$
traz um fator $2$ que a de $\lambda\abs{\beta}$ não traz; ao isolar $\beta$, dividimos
tudo por $2$, e o $\lambda$ vira $\lambda/2$.

Com o RSS definido com $\tfrac12$ na frente (e a penalidade sem), o problema seria
minimizar $\tfrac12\norm{\bm y-\mathbb{X}\bbeta}^2+\lambda\,\mathrm{Pen}(\bbeta)$.
Multiplicar um objetivo por uma constante positiva não muda o minimizador, e
multiplicando este por $2$ ele vira o problema da aula com $2\lambda$ no lugar de
$\lambda$. Basta fazer essa troca nas fórmulas:
\[
  \widehat\beta_j^{\,\text{ridge}}=\frac{\widehat\beta_j}{1+2\lambda},
  \qquad
  \widehat\beta_j^{\,\text{lasso}}
  =\operatorname{sinal}(\widehat\beta_j)\big(\abs{\widehat\beta_j}-\lambda\big)_+ .
\]
(Pôr o $\tfrac12$ também na penalidade não mudaria nada: seria multiplicar o objetivo
inteiro por uma constante.)

A lição prática: \textbf{o valor numérico de $\lambda$ depende da convenção}, não só
dos dados. Antes de comparar o $\lambda$ das notas com o de um livro, de um artigo ou
de uma biblioteca, confira qual objetivo cada um minimiza. O \texttt{Lasso} do
\texttt{scikit-learn} minimiza $\tfrac1{2n}\norm{\bm y-\mathbb{X}\bbeta}^2+\alpha\norm{\bbeta}_1$;
multiplicando por $2n$ vê-se que o seu $\alpha$ corresponde a $\lambda=2n\alpha$ na
convenção da aula. O $n$ vem de usar a média em vez da soma; o $2$ é o deste item.
\end{solucao}
